\exercisegroup

\begin{enumerate}
\def\labelenumi{\arabic{enumi})}
\tightlist
\item
  Let \(\mathbf{E}\) be a left topological vector space over a non-discrete topological field \(\mathbf{K}\) . A subset \(\mathbf{B}\) of \(\mathbf{E}\) is said to be bounded if for every neighbourhood \(\mathbf{V}\) of 0 in \(\mathbf{E}\) there exists \(\lambda \neq 0\) in \(\mathbf{K}\) such that \(\lambda \mathbf{B} \subset \mathbf{V}\) .
\end{enumerate}

\begin{enumerate}
\def\labelenumi{\alph{enumi})}
\item
  Show that if B is bounded in E, then for every neighbourhood V of 0 in E, there exists a neighbourhood U of 0 in K such that U.B ⊂ V.
\item
  Show that the closure of a bounded set in E is bounded. The union of two bounded sets is bounded. Every precompact set in E is bounded. Extend the corollaries of III, p.~4, prop. 4 to topological vector spaces over K.
\item
  Prove that if A is a bounded set in K (considered as a vector space on the left over itself) and B is a bounded set in E, then A.B is bounded in E.
\item
  Extend prop. 3 of III, p.~4 to the case where K is a metrizable topological division ring. e) Extend the notion of a quasi-complete space and its properties to topological vector spaces.
\end{enumerate}

\begin{enumerate}
\def\labelenumi{\arabic{enumi})}
\setcounter{enumi}{1}
\tightlist
\item
  \begin{enumerate}
  \def\labelenumii{\alph{enumii})}
  \tightlist
  \item
    Let E be a left topological vector space over a non-discrete topological field K. Prove that if there exists a neighbourhood V of 0 in E which is bounded (exerc. 1), than the sets \(\lambda\) V, for \(\lambda \in \mathbf{K}\) and \(\lambda \neq 0\) , form a fundamental system of neighbourhoods of 0 in E. If K is metrizable, the Hausdorff topology associated with the topology of E (GT, III, § 2, No.~6) is metrizable. If \(\mathbf{K} = \mathbf{R}\) or \(\mathbf{K} = \mathbf{C}\) , the locally convex topology on E which is the finest of the topologies coarser than the given topology on E (II, p.~80, exerc. 23) can be defined by a single semi-norm.
  \end{enumerate}
\end{enumerate}

\begin{enumerate}
\def\labelenumi{\alph{enumi})}
\setcounter{enumi}{1}
\item
  Prove that the topology of an infinite product of locally convex Hausdorff spaces (of which none is just 0) cannot be defined by a single semi-norm.
\item
  Let E be a locally convex space whose topology is defined by an increasing sequence \((p_n)\) of semi-norms. In order that the topology of E be defined by a single semi-norm, it is necessary and sufficient that there exists an integer \(n_0\) such that for every \(n \geqslant n_0\) there exists a number \(k_n \geqslant 0\) such that \(p_n(x) \leqslant k_n p_{n_0}(x)\) for all \(x \in E\) .
\end{enumerate}

\(d\) ) Let E be the vector space over \(\mathbf{R}\) consisting of infinitely differentiable numerical functions on the interval \(\mathrm{I} = [0,1]\) . For every integer \(n\geqslant 0\) , let \(p_n(f) = \sup_{0\leqslant k\leqslant n}\left(\sup_{x\in I}|f^{(k)}(x)|\right)\) (with

\(f^{(0)} = f)\) ; show that the \(p_n\) are norms on E and that the topology defined by the sequence of norms \(p_n\) cannot be defined by a single norm.

\begin{enumerate}
\def\labelenumi{\arabic{enumi})}
\setcounter{enumi}{2}
\item
  Let E be a metrizable vector space over R, d a translation invariant distance, compatible with the topology of E. Let \(|x| = d(x, 0)\) (I, p.~16). Prove that for every integer n \textgreater{} 0, we have \(\frac{1}{n}|x| \leqslant \left|\frac{x}{n}\right|\) . Deduce from this that if B is a bounded subset of E, then \(\sup_{x \in B} |x| < +\infty\) (in other words, B is bounded for the distance d (GT, IX, § 2, No.~3)). Give an example of a metrizable vector space E and an unbounded set in E which is bounded for the distance d (exerc. 2).
\item
  Let E be a topological vector space over a non-discrete metrizable topological field K. Show that if E is a Baire space and if in E there exists a countable base for the bornology formed by bounded subsets of E (III, p.~37, exerc. 1), then there exists a neighbourhood of 0 in E which is bounded, and consequently the Hausdorff topology associated with the topology of E is metrizable (III, p.~37, exerc. 2) (compare with exerc. 6).
\item
  Let \(\mathbf{E}\) be a metrizable vector space over a non-discrete valuated field \(\mathbf{K}\) . Prove that if \((\mathbf{B}_n)\) is an arbitrary sequence of bounded subsets of \(\mathbf{E}\) (III, p.~37, exerc. 1), then there exists a sequence \((\lambda_n)\) of scalars \(\neq 0\) , such that the union of the sets \(\lambda_n\mathbf{B}_n\) is bounded.
\item
  Let \(\mathrm{E}\) be a locally convex space, which is the strict inductive limit of a strictly increasing sequence \((\mathrm{E}_n)\) of locally convex Hausdorff spaces, each \(\mathrm{E}_n\) being closed in \(\mathrm{E}_{n+1}\) (II, p.~32, prop. 9).
\end{enumerate}

\begin{enumerate}
\def\labelenumi{\alph{enumi})}
\tightlist
\item
  Prove that E is not metrizable (use III, p.~5, prop. 6 and the preceding exerc. 5).\\
\item
  In order that there exist a countable base for the canonical bornology of E, it is necessary and sufficient that the canonical bornology of each \(\mathrm{E}_n\) has a countable base in \(\mathrm{E}_n\) .
\end{enumerate}

\begin{enumerate}
\def\labelenumi{\arabic{enumi})}
\setcounter{enumi}{6}
\tightlist
\item
  \begin{enumerate}
  \def\labelenumii{\alph{enumii})}
  \tightlist
  \item
    Let E be an infinite dimensional Banach space and let S be the family of compact, convex and balanced subsets of E, which is a directed set for the relation \(\subset\) . Show that E is the inductive limit of the inductive system of the Banach spaces \(E_{A}\) , where A runs through S. (Prove by contradiction that a neighbourhood V of 0 for the inductive limit topology of the topologies of \(E_{A}\) contains a ball with center 0; for this, note that otherwise, there will exist a sequence \((x_{n})\) of points of E such that \(\|x_{n}\| \leqslant 1/n^{2}\) , and which will not belong to V.) Deduce from this that there exist bounded subsets in E which are not contained in any \(E_{A}\) for \(A \subset S\) .
  \end{enumerate}
\end{enumerate}

\begin{enumerate}
\def\labelenumi{\alph{enumi})}
\setcounter{enumi}{1}
\tightlist
\item
  Let E be an infinite dimensional Banach space. On the vector space \(\prod_{m=1}^{\infty}E_{m}\) , where \(E_{m}=E\) for each m, let \(T_{n}\) denote the topology obtained by taking the product of the Banach space topology on each \(E_{m}\) for \(m\leqslant n\) , and for m\textgreater n, the finest locally convex topology on \(E_{m}\) ; \(F_{n}\) denotes the locally convex space \(\prod_{m=1}^{\infty}E_{m}\) with the topology \(T_{n}\) . Every identity map \(F_{n}\to F_{n+1}\) is continuous; show that the inductive limit space of the inductive system \(F_{n}\) is the space F obtained by assigning to \(\prod_{m=1}^{\infty}E_{m}\) the topology which is the product of the Banach space topologies on each of the factors. Deduce from this that there are bounded subsets in F which are not bounded in any \(F_{n}\) .
\end{enumerate}

\begin{enumerate}
\def\labelenumi{\arabic{enumi})}
\setcounter{enumi}{7}
\item
  Prove that in a space which is an infinite product of topological vector spaces (over \(\mathbf{R}\) or \(\mathbf{C}\) ) none just the point 0, there does not exist a countable base for the canonical bornology (first show that it is enough to prove this for the space \(\mathbf{R}^{\mathbf{N}}\) , and then use III, p.~38, exerc. 4).
\item
  Let E be the vector space over R consisting of all regulated functions on the interval I = \{0, 1\} (FVR, I, p.~4). For every integer n \textgreater{} 0, let \(V_{n}\) be the set of all functions \(f \in E\) such that \(\int_{0}^{1} \sqrt{|f(t)|} dt \leqslant 1/n\) . Show that the sets \(V_{n}\) form a fundamental system of neighbourhoods of 0 for a metrizable topology which is compatible with the vector space structure of E and that for this topology the sets \(V_{n}\) are bounded; but the convex envelope of each \(V_{n}\) is the entire space E. (Observe that every function \(f \in E\) can be written as \(f = \frac{1}{2}(g + h)\) , where g and h belong to E, and
\end{enumerate}

\[
\int_ {0} ^ {1} \sqrt {| g (t) |} d t = \int_ {0} ^ {1} \sqrt {| h (t) |} d t = \frac {1}{\sqrt {2}} \int_ {0} ^ {1} \sqrt {| f (t) |} d t.
\]

Deduce from this that the only locally convex topology coarser than the topology of E is the coarsest topology on E.

\begin{enumerate}
\def\labelenumi{\arabic{enumi})}
\setcounter{enumi}{9}
\tightlist
\item
  Let \((\mathrm{E}_{\iota})_{\iota \in I}\) be an infinite family of Hausdorff topological vector spaces, none just the point 0, over a non-discrete topological field K. Let E be the direct sum vector space of the \(\mathbf{E}_{\iota}\) , and \(\mathcal{T}_0\) the topology on E defined in I, p.~24, exerc. 14. Then a subset B of E is bounded for \(\mathcal{T}_0\) (III, p.~37, exerc. 1) if and only if B is contained in a product subspace \(\prod_{\iota \in H} \mathbf{E}_{\iota}\) , where H is a finite
\end{enumerate}

subset of I and the projections of B on each \(E_{\iota}\) for \(\iota \in H\) are bounded. Deduce (for K = R or C) that, if each \(E_{\iota}\) is a quasi-complete space, then E with \(T_{0}\) is quasi-complete.

\begin{enumerate}
\def\labelenumi{\arabic{enumi})}
\setcounter{enumi}{10}
\tightlist
\item
  Let \(\mathbf{E}\) be a topological vector space over a field \(\mathbf{K}\) with a non discrete valuation.
\end{enumerate}

\begin{enumerate}
\def\labelenumi{\alph{enumi})}
\item
  For a balanced subset A of E to absorb every bounded subset (III, p.~37, exerc. 1) of E, it is sufficient that A absorbs the set of points of every sequence \((x_{n})\) tending to 0 in E. Then A is said to be bornivorous.
\item
  Let \(u\) be a linear mapping from \(E\) into a topological vector space \(F\) over \(K\) . The image of every bounded subset of \(E\) under \(u\) is bounded in \(F\) if and only if for every sequence \((x_{n})\) of points of \(E\) tending to 0, the sequence \((u(x_{n}))\) is bounded in \(F\) .
\item
  Suppose \(\mathbf{E}\) is metrizable. Show that every bornivorous subset of \(\mathbf{E}\) is a neighbourhood of 0 in \(\mathbf{E}\) . Deduce that if \(u\) is a linear mapping from \(\mathbf{E}\) into a topological vector space \(\mathbf{F}\) over \(\mathbf{K}\) which transforms every sequence converging to 0 in \(\mathbf{E}\) into a bounded sequence in \(\mathbf{F}\) , then \(u\) is continuous on \(\mathbf{E}\) .
\end{enumerate}

\begin{enumerate}
\def\labelenumi{\arabic{enumi})}
\setcounter{enumi}{11}
\item
  Let E be a Hausdorff topological vector space over a field K with a non-discrete valuation, and F a metrizable vector space over K. If u is a continuous linear mapping from E into F such that, for every bounded subset B of F, \(u^{-1}(\mathbf{B})\) is bounded in E, show that u is an isomorphism from E onto a subspace of F.
\item
  Let \(I\) be an infinite set, and \((E_{\iota})_{\iota\in I}\) a family of locally convex spaces, none of which is 0. Let \(f\) be a linear mapping from \(E = \prod_{\iota\in I}E_{\iota}\) into a Banach space \(F\) . Show that if the image of every bounded subset of E under f is a bounded subset of F, then there exists a finite subset H of I such that for every \(\iota \notin H\) , the restriction of f to \(E_{\iota}\) (considered as a subspace of E) is null. (Argue by contradiction that if not, we can construct a bounded sequence \((x_{n})\) in E whose image under f is unbounded in F.)
\item
  Show that if the topology of a metrizable locally convex space E cannot be defined by a single norm, then there does not exist a countable base for the canonical bornology of E (using III, p.~38, exerc. 5, show that otherwise there will exist a bounded bornivorous set (III, p.~39, exerc. 11) in E, and complete the argument using III, p.~39, exerc. 11, c)).
\item
  In a Hausdorff topological vector space E over R, let A be a compact convex set and B a closed, convex and bounded set. Show that the convex envelope C of the union \(A \cup B\) is a closed set. (Consider a point z in the closure of C, but not in A, and reduce to the case where z = 0. Observe that there exists a neighbourhood V of 0 and a number \(\alpha < 1\) such that the relations \(0 \leqslant \lambda \leqslant 1\) , \(x \in A\) , \(y \in B\) , \(\lambda x + (1 - \lambda)y \in V\) imply that \(\lambda \leqslant \alpha\) . Next, for every neighbourhood W of 0 consider the set of triplets \((\lambda, x, y)\) such that \(\lambda x + (1 - \lambda)y \in W\) , \(0 \leqslant \lambda \leqslant 1\) , \(x \in A\) , \(y \in B\) .)
\item
  Let \(\mathbf{E}\) be a locally convex metrizable space satisfying the first axiom of countability, such that its completion \(\hat{\mathbf{E}}\) is a Fréchet space which satisfies the first axiom of countability.
\end{enumerate}

Show that every bounded subset B of \(\hat{E}\) is contained in the closure of a bounded subset of \(\hat{E}\) . (Reduce to the case where B is countable, arranged as a sequence \((x_{n})\) ; on the other hand, let \((p_{n})\) be an increasing sequence of semi-norms defining the topology of \(\hat{E}\) ; for each integer n consider a sequence \((y_{nk})_{k\geqslant1}\) of points of E, which converges to \(x_{n}\) and is such that \(p_{n}(x_{n}-y_{nk})\leqslant1\) for all \(k\geqslant1\) .)

\begin{enumerate}
\def\labelenumi{\arabic{enumi})}
\setcounter{enumi}{16}
\tightlist
\item
  Let \(\mathbf{A}\) be the set of increasing mappings \(\geqslant 1\) from \(\mathbf{N}\) into \(\mathbf{N}\) ; for every \(\alpha \in \mathbf{A}\) , let \(\mathbf{B}_{\alpha}\) denote the set of all points \(z = (z_n) \in \mathbf{R}^{\mathbf{N}}\) such that \(|z_n| \leqslant \alpha(n)\) for all \(n \in \mathbf{N}\) .
\end{enumerate}

\begin{enumerate}
\def\labelenumi{\alph{enumi})}
\item
  Show that the sets \(\mathbf{B}_{\alpha}\) form a base for the bornology of all bounded subsets of the space \(\mathbf{R}^{\mathbf{N}}\) . b) For every \(\alpha \in \mathbf{A}\) , the set \(\mathbf{RB}_{\alpha}\) is a vector subspace of \(\mathbf{R}^{\mathbf{N}}\) , distinct from \(\mathbf{R}^{\mathbf{N}}\) and dense in \(\mathbf{R}^{\mathbf{N}}\) ; hence there exists a linear form \(f_{\alpha} \neq 0\) (not continuous) on \(\mathbf{R}^{\mathbf{N}}\) such that \(f_{\alpha}(z) = 0\) for all \(z \in \mathbf{B}_{\alpha}\) .
\item
  Let \(\mathbf{E}\) be the vector space consisting of all mappings \(g: \alpha \mapsto (g_n(\alpha)) \in \mathbf{R}^{\mathbf{N}}\) from \(\mathbf{A}\) into \(\mathbf{R}^{\mathbf{N}}\) such that for all \(n \in \mathbf{N}\) , the sum \(p_n(g) = \sum_{\alpha} |g_n(\alpha)|\) is finite. Show that the \(p_n\) are semi-norms
\end{enumerate}

which define the topology of a Fréchet space on E.

\(d\) ) Let \(\mathbf{H}\) be the set of all \(h\in \mathbf{E}\) such that \(h(\alpha)\in \mathbf{RB}_2\) for all \(\alpha \in \mathbf{A}\) ; show that \(\mathbf{H}\) is an everywhere dense vector subspace of \(\mathbf{E}\) (observe that every \(h\in \mathbf{E}\) such that \(h(\alpha) = 0\) except for a finite number of values of \(\alpha \in \mathbf{A}\) belongs to \(\mathbf{H}\) ).

\begin{enumerate}
\def\labelenumi{\alph{enumi})}
\setcounter{enumi}{4}
\tightlist
\item
  Let \(\mathrm{E}_0 \subset \mathrm{E}\) be the vector subspace of \(\mathrm{E}\) consisting of all \(g \in \mathrm{E}\) such that \(\sum_{\alpha} |f_{\alpha}(g(\alpha))| < +\infty\) ;
\end{enumerate}

the mapping \(u: g \mapsto (f_{\alpha}(g(\alpha)))_{\alpha \in \mathbf{A}}\) is then a linear mapping from \(\mathrm{E}_0\) into the Banach space \(\mathrm{F} = \ell^1(\mathrm{A})\) (I, p.~4). Prove that \(u(\mathrm{E}_0)\) is everywhere dense in \(\mathrm{F}\) (observe that for every finite subset \(\mathrm{J}\) of \(\mathrm{A}\) , there exists \(g \in \mathrm{E}_0\) such that \(g(\alpha) = 0\) for all \(\alpha \in \mathrm{A} - \mathrm{J}\) and that the \(f_{\alpha}(g(\alpha))\) for \(\alpha \in \mathrm{J}\) take arbitrary values in \(\mathbf{R}\) ). Show that \(u^{-1}(0)\) is everywhere dense in \(\mathrm{E}_0\) (use \(d\) ). Finally, show that for every bounded subset \(C\) of \(E\) , there exists \(\alpha_0 \in A\) such that \(f_{\alpha_0}(g(\alpha_0)) = 0\) for all \(g \in C \cap E_0\) , and deduce that the closure of \(u(C \cap E_0)\) in \(F\) is not a neighbourhood of 0 in \(F\) .

\begin{enumerate}
\def\labelenumi{\alph{enumi})}
\setcounter{enumi}{5}
\tightlist
\item
  Let G be the graph of u in \(E_{0} \times F\) , a vector subspace of the Fréchet space \(E \times F\) . Show that G is everywhere dense in \(E \times F\) (observe that for every \(x \in E_{0}\) , \(x + u^{-1}(0)\) is dense in E). However, show that for every bounded subset M of G, the closure \(\overline{M}\) of M in \(E \times F\) does not contain the bounded set \(\{0\} \times U\) of \(E \times F\) , where U is the unit ball in F (if \(N = \text{pr}_{1}(M)\) , observe that because of e), \(\overline{u(N)}\) cannot contain U).
\end{enumerate}

\begin{enumerate}
\def\labelenumi{\arabic{enumi})}
\setcounter{enumi}{17}
\tightlist
\item
  In the Banach space \(\ell^1 (\mathbf{N})\) (I, p.~4) let \(e_m\) be the sequence \((\delta_{mn})_{n\geq 0}\) such that \(\delta_{mn} = 0\) for \(m\neq n\) and \(\delta_{nn} = 1\) . Define a continuous mapping from \(\ell^1 (\mathbf{N})\) in \(\mathbf{R}\) which transforms the bounded sequence of the \(e_n\) into a non bounded subset of \(\mathbf{R}\) (use Urysohn's th. (GT, IX, §4, No.~2, th. 2)).
\end{enumerate}

